% ssp245-review-summary.tex -- self-contained review summary for Charlie and Caspar
% Build: TEXINPUTS=/var/home/aoz/Dropbox/OZ-Labs/templates//: xelatex ssp245-review-summary.tex  (run twice)
\documentclass[11pt]{ozreport}

\renewcommand{\ozdocmark}{/var/home/aoz/Dropbox/OZ-Labs/templates/assets/oz-eye-gradient.pdf}

\title{Validation Review of the SSP2-4.5 Global\\ Temperature-Attributable Burden Run}
\author{Aaron Osgood-Zimmerman \quad\textperiodcentered\quad OZ Labs LLC}
\date{August 18, 2026}

\begin{document}
\maketitle

\begin{abstract}
We reviewed the first full production run of the global temperature-attributable
mortality pipeline (scenario SSP2-4.5, 204 countries, 27 climate models, years
2022--2050). All arithmetic and coverage checks pass. Against the published
2019 estimates of the methodology the pipeline implements (Burkart et al.,
\textit{The Lancet} 2021), our 2022 global total agrees to 0.1\%, country
ranks correlate at 0.993, and 159 of 204 countries fall inside the published
95\% uncertainty intervals. An apparent large shortfall in hot countries
relative to the \emph{current} GBD round turned out to be almost entirely
upward revision by GBD itself between rounds, not a pipeline defect. The main
genuine limitations we identified and quantified are: risk curves that are
undefined above a daily mean of 32.1\,\textcelsius{} (a growing share of
future exposure), a cluster of tropical and small-island countries where our
burden runs low and appears tied to how daily temperature variability is
represented, and a years-of-life-lost definition that differs from GBD's.
Recommendations are at the end; none blocks starting the remaining three
scenarios.
\end{abstract}

\section{Background and scope}

The project estimates deaths and years of life lost (YLLs) attributable to
non-optimal ambient temperature, for 204 countries, projected to 2050 (and
later 2100) under emissions scenarios, for the World Bank. The methodology
follows Burkart et al.\ (2021), the study underlying the Global Burden of
Disease (GBD) treatment of temperature: for each of 17 temperature-related
causes of death, a set of exposure-response functions (ERFs; relative-risk
curves relating daily mean temperature to mortality) is combined with daily
temperature fields, a location-specific theoretical minimum risk exposure
level (TMREL; the death-weighted optimal temperature), and cause-specific
mortality, to produce population attributable fractions (PAFs; the share of a
cause's deaths attributable to non-optimal temperature) and attributable
deaths and YLLs. We use the ERF curves, TMRELs, and code released with the
2021 paper; mortality comes from Institute for Health Metrics and Evaluation
(IHME) national forecasts; daily temperature comes from the World Bank
Climate Change Knowledge Portal (CCKP) 0.25-degree daily product.

The run under review is the reference-scenario (SSP2-RCP4.5, ``ssp245'')
production pass executed by Caspar in July--August 2026: 204 countries
$\times$ 27 climate models $\times$ 29 years (2022--2050), 500 draws per
combination, 159{,}732 output files, summarized into review tables on
August 7, 2026. This document reports the review of that run, conducted
August 14--17, 2026. Unless stated otherwise, ``our'' numbers below are
ensemble means across the 27 models of per-model draw means, for 2022.

\section{Data integrity}

The run's quality gates all pass: the additive identities (heat plus cold
equals non-optimal, for PAFs and deaths) hold to floating-point precision;
$|\mathrm{PAF}| \le 1$ everywhere (maximum 0.777); all values are finite; no
combination has an all-zero burden; and coverage is complete (every country
has all 27 models and all 29 years; two further models publish no ssp245
data, which is expected).

The initial pass surfaced three defective combinations, all traced to
operational causes and all recomputed by Caspar on August 17, 2026, then
re-verified: Colombia and Equatorial Guinea under one model in 2022 carried a
single cause (leftovers of an early benchmark run that the resume logic then
skipped), and one Vietnam combination lacked its YLL file (an interrupted
write). After recomputation every one of the 159{,}732 combinations carries
all 17 causes and YLLs.

\section{Validation against published estimates}

\subsection{The nine benchmark countries}

Table~1 of Burkart et al.\ (2021) publishes attributable deaths for nine
countries for 2019. Table~\ref{tab:nine} compares our 2022 ensemble against
it. Eight of nine of our point estimates fall inside the published 95\%
uncertainty intervals (UIs); the exception is Brazil (ratio 0.83), which is
low almost entirely on the heat side. The rank ordering across countries
matches (Spearman correlation 0.983). Note the comparison is not
year-matched: our 2022 (modeled climate, forecast mortality) against their
2019 (observed weather, GBD mortality); GBD's own numbers move little between
2019 and 2022, so the comparison is informative but should not be read to the
digit.

\begin{table}[b]
\centering
\small
\begin{tabular}{lrrr}
\hline
country & published 2019 [95\% UI] & ours 2022 & ratio \\
\hline
China          & 499{,}605 [425{,}094--573{,}429] & 555{,}481 & 1.11 \\
United States  & 110{,}588 [98{,}410--119{,}580]  & 103{,}964 & 0.94 \\
Mexico         & 19{,}055 [16{,}111--22{,}083]    & 20{,}787  & 1.09 \\
Brazil         & 17{,}286 [15{,}386--19{,}185]    & 14{,}383  & 0.83 \\
South Africa   & 8{,}815 [7{,}846--9{,}823]       & 9{,}575   & 1.09 \\
Colombia       & 4{,}820 [3{,}669--6{,}148]       & 4{,}746   & 0.98 \\
Chile          & 3{,}759 [3{,}244--4{,}267]       & 3{,}747   & 1.00 \\
Guatemala      & 1{,}258 [976--1{,}637]           & 1{,}095   & 0.87 \\
New Zealand    & 1{,}191 [996--1{,}361]           & 1{,}177   & 0.99 \\
\hline
\end{tabular}
\caption{Attributable non-optimal-temperature deaths per year: Burkart et
al.\ (2021) Table~1 values for 2019, with 95\% uncertainty intervals, against
our ssp245 ensemble means for 2022. ``Ratio'' is ours divided by published.}
\label{tab:nine}
\end{table}

\subsection{All 204 countries against the 2021 paper's estimates}

The paper's companion dataset (IHME, GHDx DOI 10.6069/7QWW-PF74) carries
attributable death rates and PAFs for all 204 countries for 2019. We verified
it against the paper's own Table~1 (PAFs agree to at most 0.015 percentage
points, with identical UIs), reconstructed absolute counts from its
per-person rates (accurate to roughly $\pm$10\%, calibrated on the nine
Table~1 countries), and compared all 204 countries.

\begin{ozresult}{Headline validation result}
Against the 2021 published estimates for 2019, our 2022 global total agrees
to 0.1\% (1.684 vs.\ 1.686 million attributable deaths per year), the
country ranking correlates at 0.993 (Spearman), and 159 of 204 countries
fall inside the published 95\% UIs. The pipeline reproduces the methodology
it implements at global and, for most countries, national scale.
\end{ozresult}

The 45 countries outside the published UIs decompose as 41 low and 4 high.
Of the lows, 38 are tropical: a cluster of 18 small, mostly island countries
(uniformly at 0.25--0.65 of the published value), and 23 larger countries
(e.g.\ Indonesia at 0.56, Haiti at 0.35). Section~\ref{sec:limits} discusses
the likely mechanism. About two thirds of all 45 sit within 15\% of a UI
bound, close enough that reconstruction noise and the year mismatch
plausibly account for them.

\subsection{The current GBD round, and an apparent hot-country shortfall}

Against the \emph{current} GBD round's published 2022 estimates (GBD 2023
release), our global total is 0.874 of theirs, with strong structure
(Spearman 0.991; 203 of 204 countries agree on the sign of the net burden)
but a marked pattern: the largest gaps are all hot, mostly lower-income
countries, with ratios of 0.38--0.58 (Niger, Sudan, Iraq, India, Nigeria,
Mali, Chad, among others).

We initially treated this as a possible pipeline defect and tested four
explanations (temperature bias, observed 2022 weather, mortality inputs,
risk-curve truncation). The decisive test became possible once the 2021
companion dataset was in hand: the same countries compared against the
\emph{2021 paper's} estimates are almost all inside the published UIs
(Niger 0.84, Sudan 1.31, Kuwait 1.12, India 0.96, Iraq 0.94, Pakistan 0.90).
GBD itself revised these countries upward between rounds by factors of
1.3--2.3 (e.g.\ Sudan 2.34, Niger 2.27, India 1.67).

\begin{ozresult}{Resolution of the hot-country gap}
The shortfall relative to the current GBD round is, for these countries,
GBD's own between-round revision, not a defect in the run: the pipeline
matches the estimates of the methodology it implements. Which target the
project \emph{should} track (the implemented 2021 methodology, or the
current GBD round) is a scoping decision, discussed in
Section~\ref{sec:rec}.
\end{ozresult}

Two country-level curiosities are worth recording. Cambodia appears 73\%
high against GBD's 2022 value but matches both GBD vintages' 2019 values;
GBD's own Cambodia estimate falls 42\% between 2019 and 2022 within the
current round, for reasons we have not identified. The United Arab Emirates
is the one hot country genuinely low against both vintages (about half);
we ruled out temperature bias there directly (its CCKP temperatures match
the ERA5 reanalysis to within 0.4\,\textcelsius{} in the mean), and its
remaining gap is unexplained at the point-estimate level, though our
uncertainty interval overlaps the published one.

\section{Structural checks}

\textbf{Trends decompose sensibly once composition is removed.} Raw
attributable-death counts rise for both heat and cold in most countries
(global cold total 1.33M in 2022 to 1.99M in 2050), which looks
counterintuitive for cold under warming. Two composition effects account
for it: population aging (age-sex-standardized cold shares \emph{fall} in
140 of 204 countries while crude shares rise in 141), and the IHME mortality
forecast shifting deaths toward cold-heavy causes at high latitudes (chronic
kidney disease deaths in the high-latitude group grow from 269k to 687k by
2050). Within-cause, within-age climate signals behave as expected: the
age-standardized heat share rises in 197 of 204 countries. Any headline
trend statement should therefore be made on age-standardized attributable
\emph{fractions}, not raw counts.

\textbf{No climate model dominates the ensemble.} The largest mean absolute
robust deviation from the per-country ensemble median belongs to known
high-climate-sensitivity models (UKESM1-0-LL, HadGEM3-GC31-LL, KACE-1-0-G),
as expected, and no model exceeds five robust standard deviations in more
than 0.8\% of combinations. A handful of near-zero-burden countries show
sign disagreement across models (e.g.\ Singapore, ensemble $-96$ deaths,
model range $-171$ to $+135$); for these, sign instability rather than
spread is the meaningful flag.

\textbf{YLLs per attributable death differ from GBD by construction.} Our
global ratio is 13.9 years per attributable death versus 21.2 (current GBD
round) and 19.4 (2021 estimates). This is consistent with GBD computing
YLLs against its standard reference life table while we use national United
Nations World Population Prospects remaining life expectancies. It is a
definitional difference, not an error, but any side-by-side YLL comparison
must state it.

\section{Quantified limitations}\label{sec:limits}

\subsection{Risk curves end at 32.1\,\textcelsius{}}

The released ERF curves are estimated only over each climate zone's observed
daily-temperature range; the hottest zone's curve ends at a daily mean of
32.1\,\textcelsius, and hotter days are evaluated at the endpoint. The
released IHME code does the same, so this truncation is the published
method, not our shortcut. Its scale, measured from the CCKP daily fields
(two models): 7.9\% of world population lived in grid cells whose zone
assignment is clamped in 2022, rising to 14.9--23.5\% by 2050 under ssp245.
In the Gulf and Sahel, 25--40\% of person-days already exceed the curve
endpoint (Kuwait's exceedances average 5.5\,\textcelsius{} above it).

To bound what truncation could hide, we re-evaluated 2022 with the endpoint
log-linearly extrapolated (slope fitted to each curve's last
0.5\,\textcelsius): heat burden multiplies by 2.30 in Kuwait, 1.60 in
Niger, 1.31 in Sudan, and 1.00 in Haiti (which has no truncation). Under
that assumption, roughly half of the Gulf/Sahel gap to the current GBD
round (in log terms) closes.

\begin{ozcaveat}{Projection consequence}
Because the population share above the curve endpoints grows with warming,
projected burdens in hot countries are effectively lower bounds under the
published method: beyond the endpoint, additional warming adds no
additional risk. The extrapolation is an assumption, not evidence
(the true risk could plateau); we recommend reporting it as a labeled
sensitivity band alongside projection results, not adopting it silently.
\end{ozcaveat}

\subsection{The tropical low cluster and the exposure reference}

The temperature input deserves a precise description: the CCKP daily
product is NASA's NEX-GDDP-CMIP6, i.e.\ CMIP6 climate models bias-corrected
by quantile mapping to the GMFD observational dataset (Princeton's Global
Meteorological Forcing Dataset, reference period 1960--2014). It is neither
raw CMIP6 nor the ERA5 reanalysis that the 2021 estimates used for
exposure.

The countries where we run genuinely low against the 2021 estimates are
overwhelmingly tropical, and where we could test it (countries with local
ERA5 extracts), the CCKP product's day-to-day temperature spread runs
30--40\% below ERA5 in low-variance tropical countries (Rwanda,
Philippines) while matching well elsewhere. Since the burden depends on
the spread of daily temperature around the optimum, damped spread produces
damped burden; the 18 small-island countries, whose entire climate signal
is one or a few ocean-blended grid cells, are the extreme case and are
uniformly low.

We tested the obvious fix, re-referencing the daily series to ERA5 by
per-pixel quantile mapping, on four countries. The verdict is mixed: it
moves Rwanda onto the published value, leaves a control country unchanged,
overshoots Costa Rica, and moves the Philippines \emph{away} from the
published value (the two observational references disagree by
1.9\,\textcelsius{} in the Philippines' population-weighted mean, and mean
and variance effects interact with the local optimum). We therefore do not
recommend re-referencing as a production change on current evidence. The
practical implication is instead cautionary: at the country level, results
in the low-variance tropics are sensitive at up to the factor-of-two level
to which observational dataset defines exposure, and no published
uncertainty interval (ours or GBD's) includes that source of uncertainty.

\subsection{Other standing flags}

\emph{Sub-grid countries.} 29 countries contain at most one 0.25-degree
grid-cell centre; their across-model spread is roughly seven times the
median and their attributable share about a quarter of resolved tropical
peers'. Their numbers should be flagged in any headline table.

\emph{Fixed TMREL.} The optimal temperature is held at its 2020 value for
all projection years (no adaptation). Defensible to 2050; a strong
assumption for 2100 that should be stated, with the pipeline's annual
TMREL derivation available as a sensitivity.

\emph{Forecast mortality.} Projected burdens mix the climate signal with
IHME's demographic and epidemiological forecasts; attributable
\emph{fractions} isolate the climate signal and should carry the headline
trends.

\section{Recommendations}\label{sec:rec}

\begin{enumerate}
\item \textbf{Proceed with the remaining three scenarios} (SSP3-7.0,
  SSP5-8.5, SSP1-2.6). Nothing found here blocks them; the run is
  reproducible, complete, and validated against its target methodology.
\item \textbf{State the alignment target explicitly in all outputs}: these
  are estimates under the Burkart 2021 (GBD 2019-era) methodology. If
  alignment with the current GBD round is later desired, the clean path is
  obtaining the current-round risk curves from IHME, not tuning ours.
\item \textbf{Report projected trends as age-standardized attributable
  fractions}, with counts secondary.
\item \textbf{Attach the truncation sensitivity to projection results} in
  hot countries (a two-model local computation suffices; no full-scale
  rerun needed).
\item \textbf{Flag the sub-grid 29 and the tropical low cluster} in
  country-level tables, with the exposure-reference caveat above.
\item \textbf{State the YLL definition} (national life expectancy, not the
  GBD reference life table) wherever YLLs are shown.
\end{enumerate}

\medskip
{\raggedright
\noindent\textit{Provenance: all numbers in this document come from the
review analyses in the project repository, run August 14--17, 2026,
applied to the production summary tables of August 7, 2026 (re-issued
August 17, 2026), the GBD Results Tool exports of July 31, 2026, and the
IHME GHDx dataset DOI 10.6069/7QWW-PF74. Repository files:}
\code{ssp245-review-findings.org} \textit{and}
\code{output/review-ssp245/}\textit{.}
\par}

\end{document}
